\documentclass[11pt]{article}
\usepackage[margin=0.86in]{geometry}
\usepackage[T1]{fontenc}
\usepackage{lmodern,microtype,amsmath,amssymb,amsthm,mathtools}
\usepackage{booktabs,longtable,array,enumitem,xcolor,xurl,hyperref,fancyhdr}
\hypersetup{colorlinks=true,linkcolor=blue!45!black,urlcolor=blue!50!black,citecolor=blue!45!black,pdftitle={Adaptive aggregation for subexponential SAT},pdfauthor={Research prepared for Vladimir Reshetnikov}}
\setlength{\headheight}{14pt}
\pagestyle{fancy}\fancyhf{}\lhead{Adaptive aggregation for subexponential SAT}\rhead{October 26, 2026}\cfoot{\thepage}
\setlength{\parskip}{5pt}\setlength{\parindent}{0pt}
\setlist{itemsep=2pt,topsep=4pt}
\newtheorem{theorem}{Theorem}[section]
\newtheorem{lemma}[theorem]{Lemma}
\newtheorem{corollary}[theorem]{Corollary}
\theoremstyle{definition}\newtheorem{definition}[theorem]{Definition}
\theoremstyle{remark}\newtheorem{remark}[theorem]{Remark}
\newcommand{\poly}{\operatorname{poly}}
\newcommand{\OO}{\mathcal O}
\newcommand{\vv}{\nu_2}
\newcommand{\pc}{\operatorname{popcount}}
\newcommand{\floor}[1]{\left\lfloor#1\right\rfloor}
\newcommand{\ceil}[1]{\left\lceil#1\right\rceil}
\newcommand{\K}{\mathcal K}
\newcommand{\smallpath}[1]{{\small\path{#1}}}
\newcommand{\sourcepin}{\texttt{474ddba696a5d385c091c23fb2995f69edd13e00}}
\title{\textbf{Adaptive aggregation for subexponential SAT}\\[5pt]\large A review of fifteen refinement packets, sharper time--space bounds,\\and a measured finite-precision Cauchy-kernel improvement}
\author{Research prepared for Vladimir Reshetnikov}
\date{October 26, 2026}
\begin{document}
\maketitle
\begin{abstract}
We review the current ProveIt SAT research tree and all eight second-round packets
received through \texttt{docs/incoming}, a total of fifteen refinement packets and two primary
manuscripts. The strongest displayed general polynomial-space contraction
coefficient is $76.90899$, but combining already available ingredients gives
$66.60514$, and approaching their parameter boundary gives $66.19272$.
We use this stronger comparison baseline.

The principal new combinatorial step chooses a partition of nine groups into
three triples separately for each exact mode-weight allocation. An averaging
argument bounds the crossing-index sum, so direct enumeration improves the
fixed-partition matrix-multiplication certificate. Relative to the reviewed
contraction framework, the resulting normalized bounds are
\[
 T\le\poly(s+\mathcal L)\,
 2^{(36.75031252384+o(1))s/\sqrt{\log_2s}},\qquad
 S\le\poly(s+\mathcal L)\,
 2^{(34.02417800750+o(1))s/\sqrt{\log_2s}}.
\]
A separate recursive separator summation gives coefficient $55.70947831891$
while retaining polynomial space. A low-degree specialization of the cached
algorithm gives time coefficient $15.3919331988$. These constants are per normalized Boolean wire,
not per SAT variable. The logarithmic order of the general SAT bound is unchanged.

We also implement the existing joint Cauchy factorization directly modulo
$2^p$, using separate exact valuations and bounded unit residues. On six
paired local-kernel fixtures it is $6.22$--$58.62$ times faster than a
strengthened rational-recurrence control with the same mode cutoff. Complete
code, exact checks, raw timings, source hashes, and this LaTeX source accompany
the report. No complete SAT solver or new proof-assistant verification is claimed.
\end{abstract}
\clearpage
\begingroup
\setlength{\parskip}{0pt}
\tableofcontents
\endgroup
\clearpage

\section{Result, comparison standard, and proof boundary}
\label{sec:result}

The task is to improve on the combined frontier, rather than repeat an idea
already in another packet. Three contributions survive that comparison:
\begin{enumerate}
\item \textbf{Adaptive aggregation.} For every exact mode-weight allocation,
choose the least expensive crossing partition among the $280$ partitions of
nine groups into triples. The new bound follows from a finite averaging
identity, not a hypothetical compact tensor compiler.
\item \textbf{Polynomial-space separator summation.} Recursively split the
mode graph and, after fixing the crossing indices, sum its two sides
independently. Exact distributivity avoids jointly enumerating the internal
indices of separate branches and reduces the paid recursive work.
\item \textbf{A practical arithmetic implementation.} Evaluate the known
Cauchy features with bounded modular arithmetic, exact valuation bookkeeping,
reused denominators, and streaming. Compare against both the unchanged
archive and a stronger matched recurrence control.
\end{enumerate}

\subsection{Normalized interface}

The common source interface is a bipartite Boolean tensor network. Its
contraction is
\[
 Z=\sum_{x\in\{0,1\}^{E}}\prod_v T_v(x|_{\delta(v)}).
\]
The root has at most $s$ wires; the initial precision and truncated coefficient
axis have size $O(s)$ under the source normalization. The root maximum degree
is at most $\ceil{s^{9/10}}$, the tensor descriptions have total length
$\mathcal L$, and original tensor entries obey the source's permitted
$\exp(O(s^{9/10}))\poly(s+\mathcal L)$ evaluation cost and polynomial-space
specification. Arithmetic takes place in an unramified extension of the
two-adic integers, optionally with a truncated nilpotent coefficient axis,
modulo a prescribed power of two. All scalar and coefficient dimensions
remain polynomial in $s$.

These are the same conditions under which the compared normalized constants
are derived. An arbitrary dense tensor supplied by an arbitrary costly oracle
does not automatically satisfy them. Tensor creation precision is part of
its immutable description: later evaluation reuses the specified canonical
representatives. Canonical lifting is not assumed to be a ring homomorphism.

\begin{theorem}[Adaptive time--space certificate]\label{thm:main}
Under the normalized interface and the inherited exact identities specified
in Section~\ref{sec:dependencies}, there is an exact contraction algorithm with
\begin{align}
 T(s)&\le\poly(s+\mathcal L)
 2^{(C_{\rm ad}+o(1))s/\sqrt{\log_2s}},\label{eq:main-time}\\
 S(s)&\le\poly(s+\mathcal L)
 2^{(C_{\rm mem}+o(1))s/\sqrt{\log_2s}},\label{eq:main-space}\\
 C_{\rm ad}&=\frac{\pi\sqrt{40\log_2 9}}{\ln2}\frac{\sqrt{42}}9
 =36.75031252383623\ldots,\notag\\
 C_{\rm mem}&=\frac{\pi\sqrt{40\log_2 9}}{\ln2}\frac23
 =34.02417800749051\ldots.\notag
\end{align}
There is also a polynomial-space implementation with coefficient
$C_{\rm ps}=55.70947831890997\ldots$, given explicitly in
Section~\ref{sec:polyspace}.
\end{theorem}

The two algorithms are different resource choices. The faster certificate
stores subexponential tables and gives up polynomial space. It does not
dominate the polynomial-space branch under a fixed memory budget.
The $o(1)$ terms and eventual parameter thresholds matter: a coefficient such
as $36.7503$ is not an observed running-time constant or a finite-size
guarantee without its correction term.

\subsection{What has and has not been proved here}
\label{sec:dependencies}

The adaptive partition identity, the streamed separator algorithm, their
resource accounting, and the new schedule analysis are proved below. The new arithmetic
implementation has an explicit algebraic correctness argument and independent
exact finite checks.

The following are imported from the reviewed sources, with their hypotheses
preserved: the deformation and exact recovery identities; construction of
the child tensor networks; the two-stage asymmetric packing lemma; the sharp
Hermite precision loss; the all-level, group-preserving joint Cauchy basis;
the original entry-oracle and canonical-representative conventions; and the
reductions from SAT. This report does not independently reconstruct all of
those arguments or rebuild the upstream Lean development. Consequently the
global SAT consequences are \emph{relative refinements of that framework},
not a newly certified independent proof of the ETH claim.

Finite checks do not prove the asymptotic theorem. Conversely, a smaller
proved upper-bound coefficient does not establish a speedup over the actual
runtime of an unimplemented algorithm. The measured claim concerns the
identified feature-table kernel only.

\section{Frozen source review}
\label{sec:review}

The closing repository snapshot is \sourcepin. The initial review used
\texttt{ab2fbc9e03d5c3d0ad6dec8686960bfe498b0b06}, where seven new SAT
archives had just arrived in \texttt{3e57705cf14ee}. A closing check found one
additional matching-potential packet in \texttt{d098ba436b33} and the filing
commit \texttt{6d3adc450615}, committed October 26, 2026 at 19:03:44 Pacific
time. All eight second-round packets are now under
\texttt{Computability/SubexponentialSAT/Research}; their historical arrival
archives remain identified in the ledger.

The closing review includes the additional matching report and verifies the
filed versions of the other seven against the reviewed archives. The final
corpus contains fifteen refinement packets and two primary manuscripts.
The initial seven-packet summary was therefore insufficient, and the report
uses the expanded closing corpus.

The first relocation commit is \texttt{2b925d8800cd}; the original preprint
was vendored by \texttt{3e096089911f}. Its upstream pin is
\texttt{7ef5c0cc5ca4b0e072b7a050fca09ab24bc6fa3a}. The accompanying
\smallpath{provenance/reviewed_sources.json} records the closing file paths,
Git blob hashes, SHA-256 hashes, and immutable URLs. The initial source and
archive inventory is retained separately. No repository branch or file was
submitted or changed.

\subsection{Primary manuscripts and seven earlier packets}

The OpenAI manuscript's Appendix D gives the general order
\begin{equation}
 \poly_k(n,m)2^{O_k(n\log\log n/\sqrt{\log n})},
 \label{eq:general-sat}
\end{equation}
and the fixed-density order $2^{O_{k,D}(n/\sqrt{\log n})}$, in polynomial
space~\cite{openai}. The separate Claude manuscript states an order based on
$\min\{km/\sqrt{\log m},k^2n\log\log n/\sqrt{\log n}\}$~\cite{claude}.
Its implementation and matched explicit network constant are not provided.
Numerical domination over that different architecture cannot be inferred
from a smaller normalized contraction coefficient.

\begin{longtable}{@{}>{\raggedright\arraybackslash}p{0.27\textwidth}>{\raggedright\arraybackslash}p{0.67\textwidth}@{}}
\toprule
Earlier packet & Established ingredient and remaining limitation\\\midrule
\endfirsthead
\toprule Earlier packet & Established ingredient and remaining limitation\\\midrule
\endhead
\path{cauchy_mode_improvement} & Integral factorial Cauchy features with
weights $1+b+2\vv(b!)$; slope three improves the mode-entropy coefficient.
No change in logarithmic order.\\
\path{eth_contraction_improvements} & Sharp taper/precision, leading
precision $6MD/u$ in place of $24MD/u$; explicit coefficient $648.44107$;
streamed finite-difference extraction.\\
\path{3sat_precision_refinement} & Exact recovery loss and a solved
real taper minimax problem. The $n\log\log n/\log n$ target requires a
missing structured summation algorithm.\\
\path{sat_density_improvement} & Uniform density-sensitive exponent
$O_k(n\min\{\log(2+m/n),\log\log n\}/\sqrt{\log n})$.
It improves some intermediate-density regimes.\\
\path{eth_density_refinement} & Overlapping density-sensitive bound;
precision refinements and a local Smith-spectrum analysis. Neither changes
the unrestricted denominator $\sqrt{\log n}$.\\
\path{3sat_refinement} & Scalar parity lift, reduced incidence-degree
cap, and removal of an auxiliary coefficient axis. Its second-contraction
target needs an unsupplied compact factorization.\\
\path{3sat_asymptotic_analysis} & Restricted barriers for taper and
mode retuning. Stronger targets are conditional on missing recursive or
linear-wire interfaces.\\
\bottomrule
\end{longtable}

These reports are filed under \texttt{Computability/SubexponentialSAT/Research}.
Their exact file identities are in the provenance ledger; the independent
review is included as \smallpath{audit/frontier_audit.md}.

\subsection{All eight second-round packets}

The following numeric coefficients use the common scale
$s/\sqrt{\log_2s}$ in a base-two exponent, under the general degree premise.
They must not be mixed with lower-degree specializations.

\begin{longtable}{@{}p{0.40\textwidth}r p{0.36\textwidth}@{}}
\toprule
Second-round packet & Coefficient & Main ingredient\\\midrule
\endfirsthead
\toprule Second-round packet & Coefficient & Main ingredient\\\midrule
\endhead
\path{eth_balanced_packing} & $94.74820$ & Nearly lossless balanced padding
and an additive precision guard.\\
\path{eth_high_moment} & $94.74820$ & Higher-moment labeling and an
alternative orbit-padding construction.\\
\path{eth_packet_refinement} & $95.66394$ & Improved pair loads, crossing
fraction, and a thin precision gap.\\
\path{eth_recursive_refinement} & $133.20737$ & Sharper packing and
an actual-domain Cauchy refinement.\\
\path{3sat_round2_refinement} & $84.06627$ & Matching-potential packing
and a combined two-stage schedule.\\
\path{eth_asymmetric_packet_balancing} & $76.90899$ & Different bounds
for the retained and fresh tensor sides.\\
\path{eth_matching_refinement} & $94.74820$ & Exact matching-potential
labeling and thin additive precision gaps; no new summation method.\\
\path{eth_joint_cauchy} & $319.68576$ & Full-orientation slope-four modes
at all retained levels; unchanged family multiplicity.\\
\bottomrule
\end{longtable}

All are polynomial-space certificates as displayed in the table. The last
packet also gives a cached, fixed-partition Strassen branch with time
coefficient $225.59706$ and space coefficient $196.19945$.
The asymmetric packet's $31.06413$ specialization assumes root degree
$s^{1/12}$, rather than $s^{9/10}$. The high-moment and recursive reports
also give lower-degree variants. These are retained as separate branches
in the comparison, not treated as weaker general algorithms.

\subsection{The stronger combined baseline}
\label{sec:baseline}

The eight second-round packets were each written against the older seven-packet
frontier. Combining the new archives with one another is therefore necessary.
Asymmetric packing preserves the old child graphs and gives
\[
 M_{\rm child}\le(2/r+o(1))M,\qquad
 c_i\le(2(1-1/r)+o(1))M/D_i.
\]
The joint Cauchy basis has exactly the same mode-family multiplicity and
source/target group locality. Its weights have slope four instead of three.
It can therefore be composed with the asymmetric schedule.

For the existing $r=16$ schedule, that composition gives
\begin{equation}
 C_{\rm comp}=\frac{\pi\sqrt{15/8}}{\ln2}(9+\sqrt3)
 =66.60514340761482\ldots.\label{eq:composed-baseline}
\end{equation}
Letting the parameters approach the admissible boundary slowly enough to
absorb their logarithmic buffers gives
\begin{equation}
 C_{\rm lim}=\frac{\pi\sqrt{150}}{\ln2}
 \left(1+\frac1{3\sqrt3}\right)
 =66.19272369357722\ldots.\label{eq:limiting-baseline}
\end{equation}
The complete integer boundary schedule is reproduced in the audit.
The new polynomial-space coefficient is compared against
\eqref{eq:limiting-baseline}, not just against $76.90899$. A preliminary
two-alphabet retuning also gives $66.12053$; its complete schedule and
independent review appear in \smallpath{audit/two_alphabet_review.md}.
The separator summation proved below improves substantially on both.

Similarly, the cached fixed-partition evaluator can be combined with the
better packing and joint modes. For instance $r=8$, block sizes $(2,3,3)$,
and geometric taper ratio $1/2$ give the limiting coefficient
\begin{equation}
 C_{\rm fixed,8}=\frac{\pi\sqrt{120}}{\ln2}
 \sqrt{\frac{30+12(\log_2 7-2)^2}{64}}
 =38.16774569701537\ldots.\label{eq:fixed-composed}
\end{equation}
Thus a comparison only with the packet's displayed $225.59706$ would greatly
overstate the contribution. The adaptive evaluator improves this stronger
composed certificate as well. We do not claim to have optimized every
possible algorithm or every unreported constant in the corpus.

\section{The shared arithmetic and recursion interfaces}
\label{sec:interfaces}

At an active call with $M$ wires, let $L=2^a$, let $r$ be the group alphabet,
and write $Q=r^a$. Label each side of the graph by $[r]^a$ and pad every
ordered label pair to a common multiple $H_{\rm pad}$ of $L$. The total padded
mass is $P=H_{\rm pad}Q^2$.

\subsection{Asymmetric padding and side provenance}

The reviewed asymmetric packing lemma gives a deterministic polynomial-time
labeling with pair load bounded by
\begin{equation}
 \left(1+\frac2{h-1}\right)
 \left(\frac M{Q^2}+\frac RQ\right)+2(hC+1)\ln Q,\label{eq:load}
\end{equation}
where $R,C$ bound row and column degrees, respectively, and $h\ge2$.
It first balances whole row vertices greedily, then derandomizes column
labels by an exact rational conditional expectation. Parallel edges are
distinct ports and are covered by the lemma.

Taking $h\asymp\log s$ and rounding upward to a packet multiple yields
$P=(1+o(1))M$ whenever, uniformly over active calls,
\begin{equation}
 \frac{QR}{M}=o(1),\qquad
 \frac{Q^2C(\log s)^2}{M}=o(1),\qquad
 \frac{LQ^2}{M}=o(1).\label{eq:padding-gates}
\end{equation}
The source orientation construction is unchanged. A child has at most
$2P/r$ wires. All retained parent tensors lie on one side of a child; fresh
packet tensors lie on the other and have degree at most $2L$. Hence below
the root we may keep
\begin{equation}
 R_d\le\Delta,\qquad C_d\le2L_{d-1},\qquad
 \Delta=\ceil{s^{9/10}},\label{eq:side-degrees}
\end{equation}
provided newly created degrees stay below $\Delta$.

\subsection{Sharp precision with an additive buffer}

For $i=\vv(t)<u\le a$, put $D_i=2^{a-i}$ and use multiplicities
$q_i=\rho(u-i)$, where
\[
 \rho=\ceil{\frac{M(2L-1)+1}{L(2u-2+2^{1-u})}}.
\]
The denominator is the exact number of Hermite conditions per unit taper.
The sharp recovery formula permits the sufficient precision
\begin{equation}
 p_t=p+6\rho D_i+\rho u(a+1).\label{eq:precision}
\end{equation}
This is an imported result; its leading constant six is already optimized
within the earlier full-loss taper model.

Let $\lambda=\ceil{\log_2s}$, stop at $M\le H=\ceil{s/\lambda}$, and set
\begin{equation}
 b=\ceil{6\log_2\lambda},\qquad
 a_d=\sum_{j=0}^d u_j+(d+1)b.\label{eq:buffer}
\end{equation}
For active depths $O(\log\lambda)$, ancestral precision obeys
$p=O(s\lambda^2 2^{a_{d-1}})$. Meanwhile
$D_i\ge2^{a_{d-1}+b+1}$. Since $s/M\le\lambda$ and $u_d\le\lambda$,
\[
 \frac{p}{MD_i/u_d}=O(\lambda^4 2^{-b})=o(1).
\]
The final summand of \eqref{eq:precision} is also lower order. When the
least active $u_d$ tends to infinity, uniformly,
\begin{equation}
 p_t\le(6+o(1))MD_i/u_d.\label{eq:precision-envelope}
\end{equation}

\subsection{The full-orientation joint Cauchy basis}

For $R=2^i$, $N=L/2$, and
\[
 s_j=(j\bmod R)+2R\floor{j/R},\qquad 0\le j<N,
\]
one crossing orientation has the scalar difference kernel
\[
 W_{xy}(T)=\frac{\Gamma(T)}{T+x-y},\quad
 \Gamma(T)=\frac{\prod_{v=-L}^{L-1}(T-v)}{\mathcal C'(0)},\quad
 T=t+2^{a+1}\zeta.
\]
Here $\mathcal C(T)=\prod_{v=-L}^{L-1}(T-v)$, $\vv(T)=i$, and
$\Gamma(T)\in2^{a+1}\OO$. The displayed quotient is interpreted after
cancelling its linear polynomial factor.

Define
\begin{align*}
 P_b(x)&=\prod_{h<b}\frac{x-s_h}{T+x-s_h},&
 Q_b(y)&=\prod_{h<b}\frac{s_h-y}{T+s_h-y},\\
 f_b(x)&=\frac{TP_b(x)}{T+x-s_b},&
 g_b(y)&=\frac{TQ_b(y)}{T+s_b-y}.
\end{align*}
The exact finite telescoping identity is
\begin{equation}
 \frac1{T+x-y}=\frac1T\sum_{b=0}^{N-1}f_b(x)g_b(y).\label{eq:cauchy}
\end{equation}
For $b=qR+r'$, $0\le r'<R$, put $E_b=q+\vv(q!)=2q-\pc(q)$.
The prior integrality lemma shows that
$F_b=2^{-E_b}f_b$ and $G_b=2^{-E_b}g_b$ are integral on the slot domain.
Thus
\begin{equation}
 W_{xy}(T)=\sum_b A_b(x)G_b(y),\qquad
 A_b(x)=\frac{\Gamma(T)}T2^{2E_b}F_b(x),\label{eq:features}
\end{equation}
with certified weight
\begin{equation}
 w_q=a+1-i+4q-2\pc(q).\label{eq:weights}
\end{equation}
There are $R$ modes of each $q$ per orientation. Each orientation's source
form reads a whole half-packet, but all those ports still belong to a single
source group; its targets belong to one other group. This is exactly the
locality needed for the unchanged child factorization.

For any ordered pair of distinct groups $g,h$, uniform label-pair padding
therefore gives at most
\begin{equation}
 c_{gh}=\frac{2P}{r^2D_i}
 =\left(\frac2{r^2}+o(1)\right)\frac M{D_i}\label{eq:pair-colors}
\end{equation}
modes at each index $q$. Summing over $r(r-1)$ pairs recovers the asymmetric
packet's total family envelope. No extra valuation-class factor is introduced.

\section{Entropy of exact pair-weight lists}
\label{sec:entropy}

The recursion sums only mode tuples whose total weighted degree is below
$p_t$. Split a tuple by the exact weights $b_{gh}$ allocated to its directed
group pairs, with $\sum_{g\ne h}b_{gh}<p_t$. For fixed $r$, there are only
polynomially many such arrays, although the polynomial exponent is large.

For a particular array, let $\mathcal I_{gh}$ be the list of mode tuples on
that pair with exact weight $b_{gh}$, and put
$h_{gh}=\log_2|\mathcal I_{gh}|$. Arrays with an empty list contribute zero
and are skipped. Each child factor $Z_g$ depends only on indices incident
with $g$. Define the convenient scale
\begin{equation}
 \K(M,u)=\frac{\pi\sqrt2}{\ln2}\frac M{\sqrt u}.\label{eq:K}
\end{equation}

\begin{lemma}[Entropy on selected directed pairs]\label{lem:entropy}
For every fixed set $I$ of directed pairs,
\begin{equation}
 \sum_{e\in I}h_e\le
 \left(\frac{\sqrt{|I|}}r+o(1)\right)\K(M,u),\label{eq:subset}
\end{equation}
uniformly over exact weight arrays and active calls in the schedules below.
More generally, fixed nonnegative weights $v_e$ give the factor
$\sqrt{\sum_ev_e^2}/r$ on the right.
\end{lemma}
\begin{proof}
Lower the actual weights to $\widehat w_q=1+4q-2\pc(q)$ and extend the
finite index range to all $q\ge0$, both valid upper bounds. The counting
function of these weights is $x/4+O(\log(x+2))$, giving
\[
 \sum_{q\ge0}-\ln(1-e^{-\theta\widehat w_q})
 =\frac{\pi^2}{24\theta}+O(\log^2(2/\theta)).
\]
This follows by integrating against the counting function and using
$\int_0^\infty-\ln(1-e^{-x})\,dx=\pi^2/6$.

Use the generating-product bound for pair $e$ with parameter
$\theta/v_e$ when $v_e>0$, multiply its logarithm by $v_e$, and sum.
With $c=(2/r^2+o(1))M/D_i$, this yields
\[
 \ln2\sum_ev_eh_e\le
 \theta p_t+\frac{c\pi^2}{24\theta}\sum_ev_e^2
 +O(c\log^2(2/\theta)).
\]
Optimizing the leading two terms and applying
\eqref{eq:precision-envelope} proves the claim. The saddle parameter is of
order $\sqrt u/D_i$. The additive buffer ensures $D_i\ge\lambda^6$ up to
a fixed factor, so the error divided by $M/\sqrt u$ tends to zero uniformly.
This joint argument also covers very small or zero individual allocated
budgets; it does not assume a separate asymptotic saddle for every pair.
\end{proof}

\subsection{Exact indexing and partition selection are computable}

For mode weights $w_1,\ldots,w_v$, the finite dynamic program
\[
 D(j,z)=D(j-1,z)+D(j,z-w_j),\quad D(0,0)=1,
\]
with other zero/negative boundary values counts exact-weight tuples.
The number of modes and $p_t$ are polynomial in $s$, and the counts have
polynomial bit length, since they are at most $(p_t+1)^v$. The table provides
exact list sizes, rank/unrank operations, and enumeration with polynomial
delay. No counting oracle is assumed.

The partition algorithm below minimizes a product of these integer list
sizes. It can compare the products exactly, without calculating floating-point
logarithms. If binary ranks are used internally, rounding each logarithm adds
at most a constant per pair and does not change the leading coefficient.

\section{The new adaptive aggregation step}
\label{sec:adaptive}

We now fix $r=9$. For each exact weight array, consider all $280$ unlabelled
partitions of the nine groups into three blocks of size three. Choose one
minimizing
\begin{equation}
 H_{\rm cross}=\sum_{g,h\text{ in different blocks}}h_{gh}.\label{eq:cut}
\end{equation}
The choice is allowed to depend on the weight array. Once chosen, it is
fixed while evaluating that array.

\begin{lemma}[Adaptive crossing bound]\label{lem:partition}
For every nonnegative directed edge weighting $h_{gh}$, there is a partition
into triples for which
\begin{equation}
 H_{\rm cross}\le\frac34\sum_{g\ne h}h_{gh}.
 \label{eq:average-cut}
\end{equation}
For the exact mode lists, this implies
$H_{\rm cross}\le(1/\sqrt2+o(1))\K(M,u)$.
\end{lemma}
\begin{proof}
In a uniformly random partition into triples, a fixed group has two partners
among the other eight groups. Thus a fixed directed pair is separated with
probability $6/8=3/4$. By linearity of expectation, the mean crossing weight
is $3/4$ of the total; its minimum is no larger. The total has $72$ directed
pairs, so Lemma~\ref{lem:entropy} bounds it by
$(\sqrt{72}/9+o(1))\K$. Multiplication by $3/4$ gives $1/\sqrt2$.
\end{proof}

\subsection{Evaluate children once, then build block tables}

For each group $g$, evaluate and store its factor $Z_g$ on all its incident
pair indices. Each evaluation is a real recursive contraction of its child;
this is not a cheap black-box oracle. A group has $16$ incident directed
pairs, so the number of recursive calls per group is at most
\begin{equation}
 2^{(\alpha+o(1))\K(M,u)},\qquad \alpha=4/9.\label{eq:alpha}
\end{equation}

For a triple $B$, build a table indexed by the directed pairs joining $B$
to its complement. Sum over its internal pair indices the product of the
three cached group factors. There are
\[
 \underbrace{2\cdot9\cdot3-3^2-3}_{42}
 \quad\text{incident directed pairs},\qquad
 \underbrace{2\cdot3\cdot6}_{36}
 \quad\text{boundary directed pairs}.
\]
For every choice of the partition, Lemma~\ref{lem:entropy} gives the bounds
\begin{align}
 \text{block construction work}&\le
 2^{(\beta+o(1))\K(M,u)},&\beta&=\sqrt{42}/9,\label{eq:beta}\\
 \text{block-table storage}&\le
 2^{(j+o(1))\K(M,u)},&j&=2/3.\label{eq:j}
\end{align}
There are only three tables and nine caches, absorbed into polynomial factors.

\subsection{Directly sum the crossing indices}

After the three block tables have been built, summing their product over all
crossing indices costs $2^{H_{\rm cross}}$ table operations. By
Lemma~\ref{lem:partition}, its coefficient is at most $1/\sqrt2$, strictly
less than $\beta$, since $1/2<42/81$. Thus the block construction is the
dominant local nonrecursive work. No matrix multiplication algorithm,
division by a nonunit, or unconstructed tensor compression is required.

\begin{theorem}[Local adaptive recurrence]\label{thm:local}
Up to polynomial factors per call, the adaptive evaluator satisfies
\begin{equation}
 T_d\le 9\,2^{(\alpha+o(1))\K_d}T_{d+1}
       +2^{(\beta+o(1))\K_d},\label{eq:recurrence}
\end{equation}
with $\alpha=4/9$, $\beta=\sqrt{42}/9$. Polynomial factors in the number
of recursive calls depend on $s$, independently of $\mathcal L$; a single
fixed polynomial in $s+\mathcal L$ pays for entry reading and arithmetic.
Its largest local table has
coefficient $j=2/3$ on scale $\K_d$.
\end{theorem}
\begin{proof}
The first term pays for every cached child contraction, and the second pays
for block-table construction and the final crossing summation. Each original
mode tuple belongs to one exact weight array. Group factors and block tables
only regroup finite sums and products in a commutative ring, so their
contraction equals the original recursion identity. Exact weight arrays,
interpolation nodes, and auxiliary evaluation choices contribute polynomial
factors, because $r=9$ is fixed and all displayed budgets are polynomial.
The preceding counts prove the bounds.
\end{proof}

\paragraph{Access cost.}
Choose the partition before constructing its group caches. Exponential tables
are not assumed to offer free random access. Store records
lexicographically in their pair indices. Cache layouts can be sorted into
the order needed by the chosen partition at polynomial overhead per record.
In a lexicographic enumeration, a projection onto a factor's axes advances,
repeats, or rewinds a suffix. Charge each rewind to the complete suffix
assignments just traversed; total movement is bounded by a polynomial factor
times the full enumeration. A global lexicographic enumeration of crossing
indices can use each block table in its induced order. The record lengths,
number of axes, and logarithms of record counts are polynomial in $s$.
Therefore the stated exponents apply to bit work, not only to an idealized
unit-cost table-lookup model.

\section{A uniform integer schedule and the global theorem}
\label{sec:schedule}

Let $\lambda=\ceil{\log_2s}$ and define, using integer comparisons,
\begin{equation}
 B_\lambda=\max\{j\ge0:9^{20j}\le2^\lambda\},\quad
 U=B_\lambda-\ceil{\sqrt\lambda},\quad
 u_d=\floor{U/2^d}.\label{eq:main-schedule}
\end{equation}
Use $r=9$, the buffer \eqref{eq:buffer}, and $H=\ceil{s/\lambda}$.
Handle bounded sizes or failed finite admissibility checks by direct
contraction. The proof below shows that all those checks eventually succeed.

\subsection{Padding, degree, and depth without circular assumptions}

First consider every hypothetical active depth
$d\le d_*:=\ceil{\log_4\lambda}+2$ and every $M>H$.
The geometric sum gives
\[
 a_d\le2U+O((\log\lambda)^2),\qquad
 u_d\ge\lambda^{1/2+o(1)}.
\]
Thus the floor errors are uniformly relative $o(1)$, and
$U/\lambda\to1/(20\log_2 9)$.

At the root, $R=C=\Delta$. By definition of $U$,
\[
 2a_0\log_2 9+\tfrac9{10}\log_2s-\log_2s
 \le-2\log_2 9\sqrt\lambda+O(\log\lambda).
\]
This tends to $-\infty$ fast enough to absorb the precision and padding
logarithms. Below the root, \eqref{eq:side-degrees} gives the same leading
negative gap for $Q_dR_d/M$, since $a_d\le2U+O((\log\lambda)^2)$.
The other two padding gates have leading exponent at most
\[
 \frac{1+2\log_2 9}{10\log_2 9}-1<0.
\]
All newly created degrees are at most
$s^{1/(10\log_2 9)+o(1)}<s^{9/10}$. This proves
\eqref{eq:padding-gates} uniformly over the hypothetical range.

Consequently the padded mass is $(1+o(1))M$, and eventually each child has
at most $M/4$ wires, since $2/9<1/4$. An active path has at most
$\log_4\lambda+O(1)$ steps, validating the hypothetical range and closing
the induction. Sharply,
\begin{equation}
 M_d\le s\left((2/9)(1+\delta_s)\right)^d,
 \qquad\delta_s\longrightarrow0.\label{eq:sharp-child}
\end{equation}

The negative $\Theta(\sqrt\lambda)$ slack is essential. Setting the
parameters exactly on the boundary and then adding the logarithmic buffers
would not justify nearly lossless padding. The inward retreat in
\eqref{eq:main-schedule} resolves this while preserving the limit coefficient.

\subsection{Solving the paid recursion}

Equations \eqref{eq:K}, \eqref{eq:main-schedule}, and
\eqref{eq:sharp-child} imply, uniformly,
\begin{equation}
 \K_d\le(\kappa+o(1))\gamma^d\frac{s}{\sqrt\lambda},\quad
 \kappa=\frac{\pi\sqrt{40\log_2 9}}{\ln2},\quad
 \gamma=\frac{2\sqrt2}{9}.\label{eq:decay}
\end{equation}
In fact, the packing relative error is $O(1/\lambda)$ plus terms that are
exponentially small in $\sqrt\lambda$ or a fixed power of $s$. Over
$O(\log\lambda)$ levels its product is $1+o(1)$. Equivalently one may keep
$\gamma_s=\gamma+o(1)$ until the uniformly convergent geometric series is
summed. A path that pays for child caches through
depth $d-1$ and performs block work at depth $d$ has coefficient at most
\[
 \kappa\left[\alpha\frac{1-\gamma^d}{1-\gamma}
                   +\beta\gamma^d\right].
\]
The bracket is a convex combination of $\alpha/(1-\gamma)$ and $\beta$.
The exact inequalities
\begin{equation}
 \frac{\alpha}{1-\gamma}
 =\frac4{9-2\sqrt2}<\frac23<\frac{\sqrt{42}}9=\beta
 \label{eq:dominance}
\end{equation}
therefore show that the root block term controls the entire recurrence.
This proves the time coefficient $\kappa\beta$ in
Theorem~\ref{thm:main}.

At any time, there are polynomially many recursion levels of resident caches
and tables. Their sizes add rather than multiply; deeper coefficient bounds
decrease geometrically. The largest local table is bounded by
$2^{(j+o(1))\K_0}$, so the global space coefficient is $\kappa j$.

The leaf cost is $2^{O(s/\lambda)}$. Original entry costs are
$\exp(O(s^{9/10}))$. Fresh packet entry cost remains
$\exp(O(L\log s))$, where $L=s^{1/(10\log_2 9)+o(1)}$ at most.
All are $2^{o(s/\sqrt\lambda)}$. Polynomial work at
$O(\log\lambda)$ levels contributes only $O(\lambda\log\lambda)$ to
the logarithm. This includes preprocessing, exact indexing, and sorting.
The precision and ring dimensions remain polynomial by the buffer argument.
This completes the theorem, including its bit-operation and memory accounting.

\section{Recursive separator summation in polynomial space}
\label{sec:polyspace}

The exact pair-weight allocation does more than enable caches. It turns the
retained mode sum into a Cartesian product of pair lists, allowing independent
internal sums to be computed without storing them. This uses ordinary
separator conditioning and distributivity, in the tradition of recursive
conditioning~\cite{darwiche}; the new contribution is its
explicit paid application and coefficient bound in this contraction framework.

Fix an exact weight array and a binary tree of the $r$ group vertices.
At a node containing a group set $B$, split $B$ into two nearly equal parts.
For each assignment to the directed pairs crossing that split, recursively
sum each side over its own internal pairs and multiply the two answers.
Assignments to pairs already fixed by ancestor splits remain parameters.
At a leaf $\{g\}$ evaluate the original factor $Z_g$ by a real recursive
child contraction. Accumulate one sum and product at a time.

\begin{lemma}[Exact streamed separator identity]\label{lem:stream-exact}
The algorithm evaluates the original finite mode sum using polynomial space.
If $E_\ell$ is the union of cut-pair sets on a root-to-leaf path, then the
number of expensive group-factor evaluations is at most
\begin{equation}
 \sum_{\ell=1}^r\prod_{e\in E_\ell}|\mathcal I_e|.
 \label{eq:stream-calls}
\end{equation}
\end{lemma}
\begin{proof}
For a split $B=B_1\sqcup B_2$, fixing its crossing indices separates all
remaining internal variables into two disjoint sets. No factor in $B_1$
depends on an internal index of $B_2$, or conversely. Finite distributivity
therefore gives a product of the two conditional sums, followed by summation
over the crossing indices. Induction on the tree proves exactness.

A leaf is visited once for each assignment to the successive cut sets on its
path. Those cut sets are disjoint, and their Cartesian product has the size
in \eqref{eq:stream-calls}. Internal indices of sibling branches are summed
independently; their domain sizes add to the work rather than multiply
through the leaf. Each leaf evaluation pays for $Z_g$ explicitly. No free
group oracle is assumed. All internal-node overhead is bounded by a fixed
polynomial times this call count. The depth is $O(\log r)$, and the assignment
ranks and ring accumulators have polynomial bit size, proving the space bound.
\end{proof}

\begin{lemma}[Balanced-tree cut profile]\label{lem:stream-profile}
For a balanced binary split on $r$ groups, the largest number of directed
pairs on a leaf path is $v(r)$, where
\begin{equation}
 v(1)=0,\qquad
 v(r)=2\floor{r/2}\ceil{r/2}
       +\max\{v(\floor{r/2}),v(\ceil{r/2})\}.
 \label{eq:v-rec}
\end{equation}
For $r=2^j$, $v(r)=2(r^2-1)/3$. In particular $v(16)=170$, whereas the
full directed group graph has $240$ pairs.
\end{lemma}
\begin{proof}
The root split contributes $2\floor{r/2}\ceil{r/2}$ directed pairs to every
leaf path. A path then enters one child, proving the recurrence. For powers
of two it becomes $v(r)=r^2/2+v(r/2)$; summing the geometric series gives
$2(r^2-1)/3$.
\end{proof}

Applying Lemma~\ref{lem:entropy} separately to each of the fixed $r$ path
sets in \eqref{eq:stream-calls} gives the local recurrence
\begin{equation}
 T_d\le\poly(s)\,
 2^{(\beta_{\rm str}(r)+o(1))\K_d}T_{d+1},\qquad
 \beta_{\rm str}(r)=\frac{\sqrt{v(r)}}r.
 \label{eq:stream-rec}
\end{equation}
A single common fixed polynomial in $s+\mathcal L$ pays for entry reading
and arithmetic. This is a reduction in the number of paid child evaluations,
not an assumption that the mode graph can be solved in subexponential time
in its index length. It still has exponential work on that length.

\begin{theorem}[Polynomial-space coefficient below $56$]\label{thm:ps}
Under the general normalized interface, streamed balanced binary summation
with $r=16$ gives
\begin{equation}
 T\le\poly(s+\mathcal L)
 2^{(C_{\rm ps}+o(1))s/\sqrt{\log_2s}},\qquad
 C_{\rm ps}=\frac{5\pi\sqrt{17}}{2\ln2}
       \left(1+\frac1{3\sqrt3}\right)
 =55.70947831890997\ldots,
 \label{eq:ps-constant}
\end{equation}
using polynomial space.
\end{theorem}
\begin{proof}
Set $g=\ceil{\sqrt\lambda}$,
\[
 B=\floor{\frac{\lambda-g}{80}},\qquad
 u_0=B,\qquad u_d=\floor{3B/4^d}\quad(d\ge1),
\]
and keep the same buffer $a_d=\sum_{j\le d}u_j+(d+1)b$ and leaf threshold.
The total retained range is at most $2B$.

To verify applicability, first consider hypothetical active depths
$d\le\ceil{\log_6\lambda}+2$. Uniformly there,
$u_d\ge\lambda^{1-\log_6 4+o(1)}\to\infty$.
At the root the $Q^2\Delta$ exponent is at most
$\lambda-g/10+O(\log\lambda)$. Below it the retained-row $Q\Delta$
exponent is at most $\lambda-g/10+O((\log\lambda)^2)$.
Fresh-side and rounding exponents are at most
$9(\lambda-g)/40+o(\lambda)<\lambda$. All padding gates hold; fresh
packet degrees remain below $\Delta$. Nearly lossless padding gives eventual
child ratio at most $1/6$, closing the depth bootstrap, and sharp ratio
$(1+o(1))/8$.

The sum of $M_d/\sqrt{u_d}$ along a path is bounded by
\[
 \left[\sqrt{80}\left(1+\frac1{3\sqrt3}\right)+o(1)\right]
 \frac{s}{\sqrt\lambda}.
\]
Multiply by $(\pi\sqrt2/\ln2)\beta_{\rm str}(16)$ with
$\beta_{\rm str}(16)=\sqrt{170}/16$. This gives
\eqref{eq:ps-constant}. Entry, precision, interpolation, and leaf costs
are lower order as in the preceding theorem. The separator recursion,
exact pair indexing, and network recursion are all streamed depth first,
so only polynomial space is retained.
\end{proof}

Relative to the combined $r=16$ polynomial-space boundary
certificate \eqref{eq:limiting-baseline}, the coefficient is multiplied by
\[
 \sqrt{\frac{170}{240}}=\sqrt{\frac{17}{24}}
 =0.8416254115\ldots.
\]
This is a $15.84\%$ reduction in that leading upper-bound coefficient. It
is an exponential separation between the displayed certificate functions,
not a measured SAT speedup. No claim of global optimality of the separator
shape or alphabet is needed. The earlier two-alphabet retuning is also
exceeded by this simpler $r=16$ construction.

\section{A lower-degree cached specialization}
\label{sec:lowdegree}

For bounded-degree clause networks it is useful to keep the root's stronger
structural premise instead of replacing it by $s^{9/10}$.

\begin{corollary}[Lower-degree time--space bound]\label{cor:low}
If the root maximum degree is at most $\ceil{s^{1/12}}$, the adaptive cached
algorithm has time and space coefficients
\begin{align}
 C_{\rm low}&=\frac{\pi}{\ln2}
  \sqrt{\frac{100\log_2(162)}{33}}\frac{\sqrt{42}}9
  =15.3919331988\ldots,\\
 C_{\rm low,mem}&=\frac{\pi}{\ln2}
  \sqrt{\frac{100\log_2(162)}{33}}\frac23
  =14.2501611298\ldots.
\end{align}
The coefficients remain per normalized wire. The original permitted tensor
entry-cost premise is unchanged.
\end{corollary}
\begin{proof}
Let $q=17/50$ and define
\[
 B_\lambda=\max\{j\ge0:162^j\le2^\lambda\},\quad
 W=B_\lambda-\ceil{\sqrt\lambda},\quad
 u_d=\floor{\frac{33}{50}W\left(\frac{17}{50}\right)^d}.
\]
The total leading retained budget is at most $W$, with limiting fraction
$A=1/\log_2(162)$. The maximum old degree at depth $d$ is bounded by
$\max\{\ceil{s^{1/12}},2L_{d-1}\}$: packet lengths increase with depth,
and retained tensors do not gain ports.

If the original-degree term dominates, the stronger symmetric packing
constraint has exponent at most $2\log_2(9)A+1/12<1$, because $1/12<A$.
If the created-degree term dominates, its exponent is at most
$(1+2\log_2 9)A=1$, with the required negative
$\Theta(\sqrt\lambda)$ slack. The rounding gate has the same bound since
$1+2\log_2 9=\log_2(162)$. All additive buffers are absorbed.

Hypothetical depth $\log_4\lambda+O(1)$ gives the least taper
$\lambda^{1+\log_4(17/50)+o(1)}\to\infty$. Padding then gives child
ratio at most $1/4$ eventually, closing the induction, and sharp ratio $2/9$.
The $\K_d$ ratio is now
$\gamma=(2/9)\sqrt{50/17}$ and its root coefficient is
$\kappa_{\rm low}=(\pi/\ln2)\sqrt{100\log_2(162)/33}$.
The exact squared inequalities
\[
 \gamma^2=\frac{200}{1377}<\left(\frac{31}{81}\right)^2,
 \qquad \left(\frac{18}{25}\right)^2<\frac{42}{81}
\]
show that $\alpha/(1-\gamma)<18/25<\beta$. The same paid recurrence
therefore has leading time $\kappa_{\rm low}\beta$ and space
$\kappa_{\rm low}(2/3)$. All other work stays lower order.
\end{proof}

This is an additional, narrower branch. It must not be advertised as the
coefficient for arbitrary $s^{9/10}$-degree inputs. Ordinary bounded-degree
SAT encodings satisfy its graph-degree premise for sufficiently large sizes;
translating the numerical coefficient to SAT variables still requires the
encoding's wire constant.

\section{A finite-precision implementation of the known Cauchy basis}
\label{sec:kernel}

The practical implementation targets exactly the feature tables in
\eqref{eq:features}. It supports scalar integer parameters and the quadratic
unramified extension
\[
 R_p=(\mathbb Z/2^p\mathbb Z)[U]/(U^2+U+1).
\]
An input pair $(u,v)$ represents the exact integer-extension parameter
$\zeta=u+vU$. This is a tested concrete specialization, not an implementation
of every extension degree required by the full contraction framework.

The recurrence for the products $P_b,Q_b$ is already present in the earlier
packet's independent verifier. We do not claim to invent that recurrence.
The new contribution is evaluating it safely at finite precision, reusing
denominators, and providing a measured implementation with a matched control.

\subsection{Why no invalid division of a residue occurs}

Each nonzero factor is stored as an exact integer valuation and a unit residue
modulo $2^p$. A slot difference never has valuation exactly $i$: it has a
lower differing bit, or is divisible by $2^{i+1}$. Hence, for nonzero
$d=s_j-s_h$,
\[
 \vv(T+d)=\min\{i,\vv(d)\}\le i.
\]
Split the numerator and denominator before modular reduction. Invert only
the unit part of the denominator. Multiplying units uses $p$-bit residues;
normalizing by $2^{E_b}$ subtracts $E_b$ from the separately tracked exact
valuation. The known integrality lemma guarantees a nonnegative remaining
valuation. No reduced residue is divided by two, and no discarded low bits
are subsequently needed.

For the prefactor $\Gamma(T)/T$, cancel the factor $T$ symbolically first.
Split every remaining factor and the scalar $\mathcal C'(0)$ denominator
into valuations and units. When $\zeta=0$, the exact zero factor is handled
explicitly. The output target/source tables are therefore exactly the prior
tables reduced modulo $2^p$.

Omit precisely the modes satisfying $w_q\ge p$, using the existing certified
weights. Every omitted target is then zero modulo $2^p$ and every source is
integral, so the reconstructed kernel is unchanged. The full-table API can
also disable truncation for verification or other callers.

\subsection{Reusing inverses and resource cost}

Denominators depend only on $d=s_j-s_b$, of which there are $O(L)$ distinct
values. Normalize and cache them once. A standard prefix-product batch
inverse computes all inverse units with one extension-ring inversion and
linear many ring multiplications. The prefactor also uses one separate
scalar modular inversion. A separate-inversion strategy is available and
is sometimes slightly faster at the small benchmark precisions.

Let $B$ be the number of retained modes, $N=L/2$, and
$S_0=NB-B(B-1)/2$. For nonzero $\zeta$, the instrumented batch variant uses
\begin{equation}
 4K+(2L-1)+5S_0-2B\label{eq:opcount}
\end{equation}
ring multiplications, where $K$ is the number of distinct cached denominators,
at most one extension-ring inversion, and one scalar inversion. When $B=0$
there is no extension-ring inversion. Thus the construction
uses $O(NB+L)$ ring operations. Generating/sorting the difference set costs
$O(NB+L\log L)$ elementary integer operations. Apart from input normalization
and integer bookkeeping, bit work is
$O((NB+L)\mathsf M(p)+\mathsf I(p))$ in this fixed-degree ring.

The table API uses $O(NB+L)$ ring storage, including its output. The streaming
API yields one pair of columns at a time using $O(B+L)$ working ring storage,
excluding consumer-retained output. Raw input normalization can involve
$a$ and the input bit length; only the accumulated unit residues are bounded
to $p$ bits. This distinction matters for arbitrary large signed parameters.

\subsection{Fair controls}

\smallpath{code/reference_joint_cauchy.py} is the unchanged archived routine.
It constructs full exact rational/extension tables with repeated products.
The stronger control, \smallpath{code/reference_recurrence.py}, adapts the
quadratic-operation recurrence already in the earlier independent verifier:
it supports both arithmetic domains, removes checking from its timed region,
and uses exactly the same mode cutoff as the new modular implementation.
Both exact controls include conversion to the requested residues. Therefore
the main comparison does not force the control to construct discarded modes
or execute its verification suite.

\section{Measured results and exact verification}
\label{sec:measurements}

The benchmark uses Python 3.12.14 on an AMD EPYC 9V74 host, one warmup and
five paired repetitions with alternating method order. All output tables
are checked for exact equality. The raw record contains every wall and
process-CPU sample, medians, parameter values, output digests, source hashes,
and operation counts. Times are local-kernel observations, not projections
from upper-bound constants.

\begin{table}[ht]
\centering\small
\begin{tabular}{@{}rrlrrrr@{}}
\toprule
$N$ & $p$ & Domain & $B$ & Exact recurrence (ms) & Modular (ms) & Gain\\\midrule
16&32&Scalar&8&2.427&0.390&$6.22\times$\\
16&32&Quadratic&8&19.039&0.524&$36.35\times$\\
32&64&Scalar&16&10.416&1.418&$7.35\times$\\
32&64&Quadratic&16&76.230&2.271&$33.57\times$\\
64&64&Scalar&32&43.966&4.541&$9.68\times$\\
64&64&Quadratic&32&391.070&6.672&$58.62\times$\\
\bottomrule
\end{tabular}
\caption{Five-run median wall times for the same finite-precision feature
tables and mode cutoff. The modular column uses the batch-inverse option.}
\label{tab:bench}
\end{table}

For the last quadratic fixture, the archived full exact implementation takes
$6.288267$ seconds, the stronger recurrence $0.391070$ seconds, and the new
modular kernel $0.006672$ seconds. The approximately $943\times$ ratio against
the archived full routine includes recurrence and truncation savings already
available in other source files; the $58.62\times$ matched-control comparison
is the more informative new result. Across the six fixtures, matched gains
range from $6.22\times$ to $58.62\times$.

These small synthetic cases were selected to exercise the arithmetic domains
and growing table sizes. They are not a SAT benchmark distribution. They do
not establish a universal speedup over every possible Cauchy implementation.
The batch inverse is not consistently faster than separately inverting the
same cached units, so the gain is attributed to bounded modular arithmetic
and reuse, with batching an optional engineering choice.

\subsection{Primary and independent exact checks}

The primary verifier passes in ordinary Python and under \texttt{python -O},
with identical canonical digests. It covers:
\begin{itemize}
\item $1{,}092$ parameter/precision cases and $35{,}868$ independent
original-polynomial matrix reconstructions;
\item $74{,}976$ target/source comparisons against the archived exact routine
and the strengthened recurrence;
\item $5{,}544$ streaming-versus-table comparisons, $1{,}092$
batch-versus-separate comparisons, and $600$ randomized unit-inversion checks;
\item eight invalid-input rejections, signed nodes, zero and nonunit
parameters, and both arithmetic domains.
\end{itemize}

An independent checker evaluates the original cancelled polynomial modulo
$2^{p+\vv(\mathcal C'(0))}$ and only then removes the known scalar denominator
valuation. It shares neither the production prefactor nor the feature
recurrence. It passes ordinary and optimized Python on $63$ further cases,
through $N=128$ and $p=129$, with $493$ reconstructed entries, $493$ checks of
the public direct-entry API, and $27$ full untruncated exact-table comparisons.
It includes zero/nonunit parameters and signed $60$-bit extension inputs.

The adaptive-aggregation checker separately verifies the partition identity,
integer-cost choice, and equality of direct and cached/block contractions on
small finite mode graphs, including a ring with zero divisors. The schedule
checker verifies exact integer capacity decisions and the coefficient
inequalities. These are implementation and finite-algebra checks; the
uniform asymptotic arguments remain the proofs in this report.

The mode-graph suite includes $240$ adaptive and $300$ streamed comparisons
with literal sums, all $280$ balanced partitions, $160$ exact integer
averaging cases, and rational certificates for the safe eventual coefficients
$56$, $37/35$, and $16/15$. A separate implementation checks $40$ recursive
conditioning fixtures over $(\mathbb Z/64\mathbb Z)[V]/(V^2)$, comparing
individual leaf-call counts and live assignment contexts as well as exact
values. Its results also agree in ordinary and optimized Python.

\section{Consequences for SAT and remaining limits}
\label{sec:sat}

The general and density-sensitive SAT orders already in the repository remain
available. Combining the reviewed reductions with the refined engine preserves
\begin{equation}
 \poly_k(n,m)\,
 2^{O_k\!\left(
 \frac n{\sqrt{\log n}}
 \min\{\log(2+m/n),\log\log n\}\right)}.
 \label{eq:density}
\end{equation}
As usual, harmless positive shifts may be placed inside the logarithms for
bounded $n$. At $m=O(n)$ this is $2^{O_k(n/\sqrt{\log n})}$; for dense
instances it gives \eqref{eq:general-sat}. Remove unused variables before
applying the reductions and multiply the exact count by their power of two.
The active variable count is at most $\min\{n,km\}$ for fixed clause width.

The faster cached evaluator changes the time--space coefficient of its
normalized contraction engine. It does not prove a denominator $\log n$ in
place of $\sqrt{\log n}$. The earlier structured-summation, compact-factorization,
and degree-preserving-recursion targets still have missing algorithmic
interfaces. Counting a compact index is not the same as cheaply evaluating
the associated tensor.

Likewise, a coefficient per wire cannot be inserted directly as a coefficient
per SAT variable. The encoding's wire constant, density dependence, reduction
branches, original tensor entry costs, and polynomial prefactors all have to
be carried through. The alternative manuscript supplies a different
big-$O$ theorem with unspecified constants, so no claim of numerical
domination over that entire algorithm is justified.

For a resource-aware portfolio, retain the earlier density and lower-degree
branches, the polynomial-space algorithm of Section~\ref{sec:polyspace}, and
the new cached branch when memory allows it. Taking the best applicable
certificate preserves every reviewed established order and adds the sharper
certificates proved here. The practical kernel is usable by either branch.

\section{Reproducible deliverable and integration}
\label{sec:deliverable}

The archive contains this complete LaTeX source and compiled PDF, standalone
standard-library Python code, raw results, independent reviews, a full source
inventory, and a build/replay script. The copied reference kernel's original
hash and the repository license are included. The project does not require
network access to replay the supplied local checks.

The principal integration changes are localized:
\begin{enumerate}
\item Retain the existing exact deformation, recovery, and immutable tensor
description conventions. Use the sharp precision and full-orientation
slope-four basis already present in the reviewed packets.
\item Use asymmetric padding and preserve the separate degree bounds of
retained and fresh tensor sides. Implement integer capacity tests and the
explicit inward slack, rather than silently using a boundary parameter.
\item For the polynomial-space branch, use balanced binary separator
summation and the $r=16$ schedule in Section~\ref{sec:polyspace}. For the
faster branch, enumerate exact pair-weight arrays, compute their
list cardinalities, choose the least crossing product among the $280$
triple partitions, build group caches and three block tables, and sum all
crossing indices. Pay for each recursive group contraction as in
\eqref{eq:recurrence}.
\item Substitute the modular feature-table implementation where its supported
coefficient ring applies. Its exact interface and streaming option are
documented in \smallpath{audit/PRACTICAL_NOTES.md}.
\end{enumerate}

The archive provides a real local arithmetic optimization and a new aggregation
prototype. It does not contain a completed end-to-end implementation of the
very large upstream SAT framework. The next meaningful engineering milestone
is an integrated contraction evaluator with memory-budget selection and
creation-precision preservation, followed by paired complete-contraction
benchmarks. Another round of abstract parameter retuning alone would not
establish a better logarithmic SAT order.

\begin{thebibliography}{99}
\small
\setlength{\itemsep}{3pt}
\raggedright
\bibitem{openai} OpenAI, \emph{The Exponential Time Hypothesis is False},
October 25, 2026. Pinned source directory:
\url{https://github.com/openai/math/tree/7ef5c0cc5ca4b0e072b7a050fca09ab24bc6fa3a/preprints/The-Exponential-Time-Hypothesis-is-False-October-25-2026}.
The reviewed files include \texttt{build/main\_full\_version.tex},
\texttt{quantitative-contraction.tex}, and \texttt{quantitative-reduction.tex}.
\bibitem{repo} Vladimir Reshetnikov and contributors, ProveIt, SAT research
snapshot \sourcepin:
\url{https://github.com/VladimirReshetnikov/ProveIt/tree/474ddba696a5d385c091c23fb2995f69edd13e00/Computability/SubexponentialSAT}.
\bibitem{incoming} ProveIt, second-round SAT packets, first arrival commit
\texttt{3e57705cf14ee9ac99c0ac35cfff751dcf70b396}:
\url{https://github.com/VladimirReshetnikov/ProveIt/commit/3e57705cf14ee9ac99c0ac35cfff751dcf70b396}.
Additional arrival \texttt{d098ba436b33}; final placement
\texttt{6d3adc45061535ebe6fa7cd3d9b37dff7dcfc22a}.
\bibitem{asymmetric} \emph{Asymmetric packet balancing for subexponential SAT},
October 26, 2026. Incoming archive
\texttt{eth\_asymmetric\_packet\_balancing\_2026-10-26.zip}; especially
\texttt{thm:pair-balance}, \texttt{cor:padding}, \texttt{lem:side-degree},
and its depth-allocation appendix. Immutable archive path and hash are in
\texttt{provenance/reviewed\_sources.json}.
\bibitem{joint} \emph{Full crossing-block Cauchy modes}, October 26, 2026.
Incoming archive \texttt{eth\_joint\_cauchy\_2026-10-26.zip}; especially
\texttt{thm:modes}, \texttt{sec:compatibility}, \texttt{lem:subsetentropy},
and the cached group/block evaluation sections. Exact identities and prototype
reference source are credited to this packet.
\bibitem{precision} \emph{Sharper precision for subexponential 3SAT},
\texttt{Research/3sat\_precision\_refinement}; exact taper, precision loss,
and full-loss minimax theorem.
\bibitem{density} \emph{Improving the clause-density dependence of exact
$k$-SAT counting}, \texttt{Research/sat\_density\_improvement}, and
\emph{A density-sensitive refinement of exact \#k-SAT counting},
\texttt{Research/eth\_density\_refinement}; both in the pinned ProveIt snapshot.
\bibitem{barriers} \emph{Subexponential 3SAT asymptotic analysis},
\texttt{Research/3sat\_asymptotic\_analysis}; report and
\texttt{research\_decision/findings.md}, pinned snapshot as above.
\bibitem{claude} \emph{A $2^{o(n)}$ time algorithm for 3-SAT}, credited to
Claude (Anthropic), October 26, 2026; filed as
\texttt{Research/claude-2o-n-3sat/eth.pdf}. Separate manuscript; no supplied
matched explicit contraction constant or public formalization build in the
reviewed packet.
\bibitem{darwiche} Adnan Darwiche, ``Recursive conditioning,''
\emph{Artificial Intelligence} 126(1--2), 5--41 (2001).
\url{https://doi.org/10.1016/S0004-3702(00)00069-2}.
The general conditioning technique is prior work; its mode-graph
application and explicit certificates are the comparison made here.
\end{thebibliography}
\end{document}
